\section{期间与送达}
\begin{enumerate}
  \item 每年考试都考。
  \item 分类。唯一依据是确定期间的方式不同。
        \begin{enumerate}
          \item 法定期间。法律明文规定。
                \begin{enumerate}
                  \item 绝对不可变。
                  \item 相对不可变。\zd{院长批准}可以延长。
                \end{enumerate}
          \item 指定期间。法院为当事人指定的期限。
        \end{enumerate}
  \item 计算方法。
        \begin{enumerate}
          \item 以时日月年为单位计算。期间开始的时和日不计算在内。当天不算，从\zd{第二天起算}。
          \item 最后一天是法定休假日顺延到第一个工作日。
          \item \zd{文书在途期间不计算在内}。
          \item 期间的耽误：障碍消除后的 10 天内\zd{申请}顺延。是否允许法院决定。
        \end{enumerate}
  \item \zd{送达}。
        \begin{enumerate}
          \item 直接送达。本人 \hz 成年家属。送达回证。\zd{离婚诉讼中只能送达本人}；通知送达；偶遇送达。
          \item 留置送达。当事人拒绝签收 \yjt 邀请基层组织工作人员见证签字 \hz 拍照录像方式。例外：\zd{调解书}。
          \item 电子送达。受送达人同意 + 确保收悉 = 电子方式送达。\zd{受送达人特定系统之日为送达之日}。
          \item 委托送达。如受送达人在外地的，\zd{其他法院}。
          \item 转交送达。他在的地方我们去不了，如军人、被监禁人等。转交对应机构、单位、监狱等。
          \item 邮寄送达。\zd{挂号信回执上记载日期为送达之日}。
          \item 公告送达。\zd{兜底送达}。国内 30 天，涉外 3 个月；不适用\zd{简易程序、支付令}（保障债务人提出异议权利）。
        \end{enumerate}
\end{enumerate}
